\documentclass[11pt,a4paper]{article}

% ---------- Encoding, page layout, and typography ----------
\usepackage[utf8]{inputenc}
\usepackage[T1]{fontenc}
\usepackage{lmodern}
\usepackage[margin=1in]{geometry}
\usepackage{microtype}
\usepackage{setspace}
\usepackage{fancyhdr}
\usepackage{titlesec}
\usepackage{enumitem}

% ---------- Mathematics, tables, and figures ----------
\usepackage{amsmath,amssymb}
\usepackage{booktabs}
\usepackage{array}
\usepackage{tabularx,longtable}
\usepackage{pdflscape}
\usepackage{listings}
\usepackage{caption}
\usepackage{placeins}
\usepackage{graphicx}
\usepackage{xcolor}
\usepackage{xurl}
\usepackage{hyperref}
\usepackage[nameinlink,noabbrev]{cleveref}

\definecolor{AccentBlue}{HTML}{1F4E79}
\definecolor{CodeGray}{HTML}{F4F7FA}
\hypersetup{
  colorlinks=true,
  linkcolor=AccentBlue,
  citecolor=AccentBlue,
  urlcolor=AccentBlue,
  pdftitle={YOLOv1: Foundations and Evolution},
  pdfauthor={Le Khoi Nguyen}
}
\setstretch{1.08}
\setlist{noitemsep,topsep=3pt,leftmargin=1.45em}
\setlength{\parindent}{1.25em}
\setlength{\parskip}{0pt}
\setlength{\headheight}{14pt}
\setcounter{tocdepth}{2}
\captionsetup{font=small,labelfont=bf,skip=6pt}
\graphicspath{{figures/}}
\newcolumntype{P}[1]{>{\raggedright\arraybackslash}p{#1}}
\lstset{basicstyle=\ttfamily\footnotesize,backgroundcolor=\color{CodeGray},
  frame=single,rulecolor=\color{AccentBlue!30},breaklines=true,
  columns=fullflexible,keepspaces=true,showstringspaces=false,
  keywordstyle=\color{AccentBlue},commentstyle=\color{gray},
  aboveskip=8pt,belowskip=8pt}
\titleformat{\section}{\large\bfseries\color{AccentBlue}}{\thesection}{0.6em}{}
\titleformat{\subsection}{\normalsize\bfseries}{\thesubsection}{0.6em}{}
\titlespacing*{\section}{0pt}{1.35em}{0.55em}
\titlespacing*{\subsection}{0pt}{1.0em}{0.35em}

\pagestyle{fancy}
\fancyhf{}
\lhead{DPL302m Personal Task}
\rhead{Le Khoi Nguyen}
\cfoot{\thepage}
\renewcommand{\headrulewidth}{0.35pt}

% Original vector schematics generated by scripts/generate_figures.py.
\newcommand{\reportfigure}[3]{%
  \begin{figure}[!htbp]
    \centering
    \includegraphics[width=\linewidth]{#1}
    \caption{#2}\label{#3}
  \end{figure}
}
% Reserved figures requested in notes/report-notes.txt.
\newcommand{\reservedfigure}[4]{%
  \begin{figure}[!htbp]
    \centering
    \fcolorbox{AccentBlue!35}{CodeGray}{%
      \parbox[c][#2][c]{0.92\linewidth}{\centering\small\color{AccentBlue}#1}%
    }
    \caption{#3}\label{#4}
  \end{figure}
}

\begin{document}

\pagenumbering{gobble}
\begin{titlepage}
  \centering
  \vspace*{1.2cm}
  {\Large\textsc{DPL302m --- Personal Task}\par}
  \vspace{1.5cm}
  {\color{AccentBlue}\rule{\linewidth}{1.2pt}\par}
  \vspace{0.55cm}
  {\LARGE\bfseries A Technical Survey of YOLOv1:\\[0.25em]
  Unified, Real-Time Object Detection\par}
  \vspace{0.45cm}
  {\large Foundations and a concise\\
  perspective on the YOLO family (v2--v12)\par}
  \vspace{0.55cm}
  {\color{AccentBlue}\rule{\linewidth}{1.2pt}\par}
  \vspace{1.7cm}
  {\large\textbf{Author}\par}
  \vspace{0.25cm}
  {\Large Le Khoi Nguyen\par}
  \vfill
  {\large October 2026\par}
\end{titlepage}
\clearpage
\pagenumbering{arabic}
\setcounter{page}{1}

\begin{abstract}
You Only Look Once (YOLO) reframed object detection as a unified regression problem: a single neural network maps an entire image directly to bounding boxes and class probabilities. This report studies the original YOLOv1 paper, published at CVPR 2016, as its primary model. It explains the problem motivation, grid-based prediction, the 24-convolutional-layer architecture, responsibility assignment, training objective, and inference pipeline. The analysis connects the model's innovations to its speed, contextual reasoning, localization errors, and difficulty with clustered small objects. A shorter historical section explains how subsequent YOLO branches address representation, training, deployment, and attention, supported by a version catalogue in the appendix. The original model occupies approximately four-fifths of the main technical discussion. The implementation section and spaces for an architecture illustration and real-image inference results are reserved for later completion.
\end{abstract}

\noindent\textbf{Keywords:} object detection; YOLOv1; one-stage detector; convolutional neural network; real-time vision

\begingroup
\small\setstretch{1.0}
\tableofcontents
\endgroup
\clearpage

\section{Introduction}
Object detection requires a model to answer two related questions: \emph{what} objects are present and \emph{where} they occur in an image. Before YOLO, many successful detectors were assembled from multiple components. A typical region-proposal pipeline generated candidate regions, extracted CNN features for each region, classified the features, refined the bounding boxes, and finally removed duplicate detections. Although accurate, this design was computationally costly and difficult to optimize as a single system \cite{redmon2016yolo,rcnn2014,fastrcnn2015}.

YOLOv1 introduced an influential alternative: predict object locations and classes from the full image in one network evaluation \cite{redmon2016yolo}. Rather than treating detection as a collection of independent subproblems, it learns a direct mapping from image pixels to spatially separated bounding boxes and associated class probabilities. The central question of this report is how this formulation achieves real-time detection, and what constraints follow from the original design.

\subsection{Scope and organization}
The report focuses on the original YOLOv1 paper and its PASCAL VOC configuration. \Cref{sec:background} explains the motivation; \cref{sec:yolov1} develops the main idea, architecture, and learning procedure; and \cref{sec:innovations,sec:assessment,sec:applications} assess innovations, strengths, limitations, and applications. \Cref{sec:demo} is reserved for an implementation example once the required materials are available. \Cref{sec:future} provides supporting historical context rather than a second full survey. Later-version references combine the supplied original papers with official software documentation for YOLOv5, YOLOv8, and YOLO11. Benchmark results in this report are attributed to the original papers; no implementation experiment is reported.

\section{Background and Motivation}
\label{sec:background}
\subsection{From multi-stage detection to unified prediction}
Sliding-window detectors repeatedly evaluate a classifier at many positions and scales. Region-based systems reduce this repetition by generating proposals, but then still process a large set of candidate regions. For example, R-CNN relies on approximately 2,000 Selective Search proposals per image \cite{rcnn2014,redmon2016yolo}. These approaches separate proposal generation, feature extraction, classification, localization, and post-processing.

YOLOv1 removes the proposal-generation and per-region prediction stages. It processes the entire image at once and learns the following direct prediction task:
\begin{equation}
  \text{image pixels} \longrightarrow \text{bounding boxes + confidence scores + class probabilities}.
\end{equation}
This formulation permits end-to-end optimization of the main prediction network and makes real-time detection a central design goal.

\reportfigure{pipeline_comparison.pdf}{A proposal-based R-CNN-style workflow and YOLOv1's unified prediction workflow. Original schematic based on \cite{rcnn2014,redmon2016yolo}; both retain post-processing.}{fig:pipelines}

The comparison in \cref{fig:pipelines} describes R-CNN, not every later region-based detector. Fast R-CNN shares a convolutional feature map across proposals and is therefore more efficient than running a CNN separately on each crop \cite{fastrcnn2015}. YOLO's distinction remains its dense, direct full-image prediction without an external proposal stage. Likewise, ``end-to-end'' in the original paper concerns joint learning of the prediction network; NMS remains outside that learned mapping.

\subsection{Why full-image reasoning matters}
Unlike a detector that classifies isolated crops, YOLO sees the complete image during both training and inference. Full-image features provide context about object appearance, scene layout, and relationships between objects. The original paper reports that YOLO produces fewer false positives on background than Fast R-CNN, though it makes more localization errors \cite{redmon2016yolo}.

\section{YOLOv1: Unified Real-Time Object Detection}
\label{sec:yolov1}
\subsection{Detection formulation and grid responsibility}
YOLOv1 conceptually divides the image into an $S\times S$ grid. A grid cell is responsible for an object if the object's center lies inside that cell. Importantly, the cell determines responsibility only; the predicted bounding box may extend beyond the cell boundary.

For the PASCAL VOC experiments in the original paper, the configuration is
\begin{equation}
  S=7, \qquad B=2, \qquad C=20,
\end{equation}
where $B$ is the number of bounding-box predictors per cell and $C$ is the number of classes. The network therefore returns a tensor of shape
\begin{equation}
  S\times S\times(B\cdot5+C)=7\times7\times30.
\end{equation}

\subsection{Output representation}
Each cell predicts $B=2$ bounding boxes. A box has five values,
\begin{equation}
  (x,y,w,h,\text{confidence}),
\end{equation}
and the cell also predicts one shared vector of $C=20$ conditional class probabilities, $P(\text{Class}_i\mid\text{Object})$. Thus, one cell produces $2\times5+20=30$ values.

The coordinates $(x,y)$ describe the box center as offsets relative to the responsible grid cell. Width and height are normalized by the image dimensions:
\begin{equation}
  w=\frac{W_{\text{box}}}{W_{\text{image}}}, \qquad
  h=\frac{H_{\text{box}}}{H_{\text{image}}}.
\end{equation}
The confidence target is defined as
\begin{equation}
  \text{confidence}=P(\text{Object})\times\operatorname{IoU}_{\text{pred}}^{\text{truth}},
\end{equation}
where IoU denotes intersection over union. At inference time, the class-specific score for box $j$ and class $i$ is
\begin{equation}
  \text{Score}_{i,j}=P(\text{Class}_i\mid\text{Object})\,P(\text{Object})\,\operatorname{IoU}_{\text{pred}}^{\text{truth}}.
\end{equation}
The class vector is shared by the two box predictors in one cell; consequently, the two predictors are not independent class classifiers.

\reportfigure{grid_representation.pdf}{Grid responsibility and the $7\times7\times30$ representation. The marked center assigns the illustrated box to one cell; the box itself spans several cells. Original schematic based on \cite{redmon2016yolo}.}{fig:grid}

For example, let an object's center lie at image-normalized coordinates $(0.35,0.493)$. With zero-based indexing, it belongs to column $\lfloor7(0.35)\rfloor=2$ and row $\lfloor7(0.493)\rfloor=3$. The corresponding local offsets are approximately $(0.45,0.45)$. A width of $0.30$ describes 30\% of the image width, not 30\% of the cell width. This distinction is essential for correct decoding: grid-relative center offsets and image-relative sizes use different reference frames.

\subsection{Network architecture}
YOLOv1 uses 24 convolutional layers followed by two fully connected layers \cite{redmon2016yolo}. Its convolutional backbone is inspired by GoogLeNet, but it does not use Inception modules. Instead, it often uses a $1\times1$ convolutional reduction layer before a more expensive $3\times3$ convolution. This reduces channel dimensionality and computational cost while retaining learned combinations of input features.

At a high level, an input image of $448\times448\times3$ is transformed into a $7\times7\times1024$ feature representation. The fully connected layers map this representation through a 4,096-unit hidden layer to the final $7\times7\times30$ detection tensor. The progressive reduction in spatial resolution provides semantic abstraction, but it also contributes to the model's coarse localization.

\begin{table}[!htbp]
  \centering
  \caption{High-level YOLOv1 architecture for the PASCAL VOC configuration.}
  \label{tab:architecture}
  \begin{tabular}{@{}p{0.22\linewidth}p{0.45\linewidth}p{0.24\linewidth}@{}}
    \toprule
    Stage & Main role & Output \\
    \midrule
    Input & RGB image & $448\times448\times3$ \\
    Convolutional backbone & Feature extraction (24 convolutional layers) & Approx. $7\times7\times1024$ \\
    Fully connected layer & Global hidden representation & $4096$ \\
    Prediction layer & Boxes, confidences, and class probabilities & $7\times7\times30$ \\
    \bottomrule
  \end{tabular}
\end{table}

\reservedfigure{Architecture illustration to be added.}{5.3cm}{Reserved space for the YOLOv1 architecture diagram.}{fig:architecture}

The architecture separates feature extraction from global prediction. The $1\times1$ layers mix channels and reduce intermediate width, while $3\times3$ layers extract spatial patterns. Four pooling stages and a later stride-2 convolution progressively reduce resolution. Flattening the final feature representation enables fully connected prediction using global information, but ties the predictor to a fixed feature size. Thus, resizing an arbitrary image to $448\times448$ is part of this original configuration; it is not a fully convolutional, arbitrary-resolution detector.

Fast YOLO is a separate, smaller network with nine convolutional layers and two fully connected layers. Its higher speed is obtained through reduced architectural capacity, with a corresponding loss of accuracy \cite{redmon2016yolo}. It should not be described as merely a faster runtime setting of the 24-layer model.

\subsection{Training strategy}
\subsubsection{Pretraining and detection conversion}
The first 20 convolutional layers are pretrained on the ImageNet 1000-class classification task using $224\times224$ inputs. The detector is then formed by adding four convolutional layers and two randomly initialized fully connected layers. Training for detection uses $448\times448$ inputs, which supplies more spatial detail for localization \cite{redmon2016yolo}.

Hidden layers use the leaky rectified linear unit (Leaky ReLU),
\begin{equation}
  \phi(z)=
  \begin{cases}
    z, & z>0,\\
    0.1z, & z\leq0,
  \end{cases}
\end{equation}
while the final layer is linear. The linear output layer allows continuous regression to coordinates and confidence-related values; the desired ranges are encouraged by target encoding and loss optimization rather than enforced by a final bounding activation.

\subsubsection{Loss function and responsibility assignment}
YOLOv1 uses a re-weighted, conditional sum-squared-error (SSE) objective. Its logic can be summarized as
\begin{equation}
  \mathcal{L}=5\mathcal{L}_{xy}+5\mathcal{L}_{\sqrt{w},\sqrt{h}}
  +\mathcal{L}_{\text{conf,obj}}+0.5\mathcal{L}_{\text{conf,noobj}}
  +\mathcal{L}_{\text{class}}.
  \label{eq:loss-summary}
\end{equation}
The coordinate weight $\lambda_{\mathrm{coord}}=5$ increases the influence of localization error. The no-object confidence weight $\lambda_{\mathrm{noobj}}=0.5$ prevents the many background predictions from overwhelming object-related gradients. Square roots of $w$ and $h$ make the objective relatively more sensitive to size errors for small boxes.

For an object-bearing cell, only the box predictor with the highest current IoU against the ground-truth box becomes responsible for coordinate and object-confidence learning. This assignment encourages the two predictors to specialize in different box shapes, sizes, or patterns. Classification loss is active only when an object is present in the cell.

For clarity, let $i$ index cells, $j$ index box predictors, and $k$ index classes. Predictions have no hat and targets have a hat. Let $\mathbb{1}^{\mathrm{obj}}_i$ denote an object-bearing cell, $\mathbb{1}^{\mathrm{obj}}_{ij}$ its responsible predictor, and $\mathbb{1}^{\mathrm{noobj}}_{ij}$ a predictor assigned a no-object confidence target. The full objective is
\begin{align}
\mathcal{L}={}&\lambda_{\mathrm{coord}}\sum_{i,j}\mathbb{1}^{\mathrm{obj}}_{ij}
 [(x_{ij}-\hat{x}_i)^2+(y_{ij}-\hat{y}_i)^2] \notag\\
&+\lambda_{\mathrm{coord}}\sum_{i,j}\mathbb{1}^{\mathrm{obj}}_{ij}
 [(\sqrt{w_{ij}}-\sqrt{\hat{w}_i})^2+(\sqrt{h_{ij}}-\sqrt{\hat{h}_i})^2] \notag\\
&+\sum_{i,j}\mathbb{1}^{\mathrm{obj}}_{ij}(C_{ij}-\hat{C}_{ij})^2
 +\lambda_{\mathrm{noobj}}\sum_{i,j}\mathbb{1}^{\mathrm{noobj}}_{ij}C_{ij}^{2} \notag\\
&+\sum_i\mathbb{1}^{\mathrm{obj}}_i\sum_{k=1}^{20}(p_i(k)-\hat{p}_i(k))^2.
\label{eq:full-loss}
\end{align}
Here, $\hat{C}_{ij}$ is the IoU target for the responsible predictor, while a no-object confidence target is zero. The notation describes positive box sizes in the paper's loss; implementations that output square-root size parameters must decode those parameters consistently. It does not imply a sigmoid-constrained final layer. The indicator structure in \cref{eq:full-loss} prevents background cells from contributing coordinate or class loss.

\reportfigure{loss_assignment.pdf}{Conditional learning signals in YOLOv1: object presence activates the class term, while maximum-IoU responsibility selects the coordinate-learning predictor. Original schematic based on the loss in \cite{redmon2016yolo}.}{fig:loss}

\subsubsection{Optimization and regularization}
The authors train for approximately 135 epochs on PASCAL VOC 2007 and 2012 training/validation data. The reported batch size is 64, momentum is 0.9, and weight decay is 0.0005. The learning rate is warmed up from $10^{-3}$ to $10^{-2}$, then maintained at $10^{-2}$ for 75 epochs, reduced to $10^{-3}$ for 30 epochs, and finally reduced to $10^{-4}$ for 30 epochs \cite{redmon2016yolo}.

To reduce overfitting, YOLOv1 applies dropout with rate 0.5 after the first fully connected layer. Data augmentation includes random scaling and translation up to 20\% of image size, together with random exposure and saturation changes in HSV space.

\subsection{Inference and post-processing}
At test time, one forward pass generates $7\times7\times2=98$ candidate boxes. Each candidate is combined with the class-probability vector from its cell to obtain class-specific scores. Low-scoring predictions are discarded, and non-maximum suppression (NMS) removes duplicate overlapping boxes of the same class. The score threshold and NMS IoU threshold are distinct hyperparameters.

\reportfigure{inference_pipeline.pdf}{YOLOv1 inference, including box decoding and post-processing. The forward network is evaluated once, but score filtering and NMS remain separate operations. Original schematic based on \cite{redmon2016yolo}.}{fig:inference}

Ground-truth boxes are not available during inference. The network predicts the confidence value that was trained toward $P(\mathrm{Object})\times\mathrm{IoU}$; it does not calculate an IoU against ground truth at test time. This distinction avoids incorrectly introducing labels into the inference pipeline. NMS instead compares pairs of predicted boxes, retaining high-scoring predictions and suppressing lower-scoring overlaps within the same class.

\FloatBarrier
\section{Key Innovations Relative to Earlier Detectors}
\label{sec:innovations}
YOLOv1's central innovation is a system-level formulation rather than a new convolution operator. Its explicit grid representation makes detection a fixed-size prediction problem, and its shared network learns localization and classification together. \Cref{tab:previous} contrasts the computation patterns of representative preceding detectors. The comparison concerns their original designs, not all modern two-stage systems.

\begin{table}[!htbp]
\centering\small
\caption{Prediction workflows before and in YOLOv1 \cite{rcnn2014,fastrcnn2015,redmon2016yolo}.}
\label{tab:previous}
\begin{tabularx}{\linewidth}{@{}lXXX@{}}
\toprule
Design & Candidate generation & Feature computation & Prediction unit \\
\midrule
R-CNN & External region proposals & Separate CNN evaluation for crops & Proposed region \\
Fast R-CNN & External region proposals & Shared feature map; RoI pooling & Proposed region \\
YOLOv1 & No external proposals & One full-image CNN evaluation & Grid cell and its two box predictors \\
\bottomrule
\end{tabularx}
\end{table}

The resulting advantages are connected: avoiding per-crop computation reduces repeated work; global features supply context; and one detection objective coordinates the learned components. However, a fixed output representation transfers complexity into the model's capacity constraints. In particular, two predictors in one cell provide alternative geometry, not two fully independent class distributions. The formulation is therefore efficient without being unrestricted.

\section{Experimental Results and Assessment}
\label{sec:assessment}
\subsection{Reported performance}
The original paper reports a clear speed--accuracy trade-off on PASCAL VOC 2007. Standard YOLO reaches 63.4\% mean average precision (mAP) at 45 frames per second (FPS), whereas Fast YOLO reaches 52.7\% mAP at 155 FPS \cite{redmon2016yolo}. These historical values should be interpreted in the hardware and benchmark context used by the original paper, rather than as direct comparisons to current detectors.

\begin{table}[!htbp]
  \centering
  \caption{Reported YOLOv1 results on PASCAL VOC 2007 \cite{redmon2016yolo}.}
  \label{tab:results}
  \begin{tabular}{@{}lcc@{}}
    \toprule
    Model & mAP (\%) & FPS \\
    \midrule
    Fast YOLO & 52.7 & 155 \\
    YOLO & 63.4 & 45 \\
    \bottomrule
  \end{tabular}
\end{table}

Mean average precision summarizes class-wise precision--recall performance. In the VOC 2007 evaluation used here, a detection must match a ground-truth object at an IoU threshold of 0.5; VOC 2007 uses its historical 11-point interpolated AP calculation \cite{everingham2010voc}. These values should not be mixed with COCO AP averaged over thresholds from 0.50 to 0.95. FPS is throughput in the original experimental setting, not a guaranteed latency on another machine.

The paper reports GPU timing on an NVIDIA Titan X and evaluates different detectors with different architectures and training histories. Accordingly, \cref{tab:results} records published evidence rather than a reproduction. A fair new comparison would hold dataset split, resolution, hardware, batch size, numerical precision, runtime, and measured pipeline stages constant. No new VOC accuracy or neural-network timing is claimed in this report.

\subsection{Advantages}
The original model offers several advantages supported by the formulation and the paper's evaluation:
\begin{itemize}
  \item \textbf{Unified prediction:} bounding boxes and class probabilities are produced by a single network evaluation.
  \item \textbf{Real-time throughput:} the detector avoids processing thousands of proposal regions independently.
  \item \textbf{Global context:} predictions use information from the complete image, reducing background false positives relative to Fast R-CNN in the original analysis.
  \item \textbf{Joint optimization:} the core detector can be optimized end-to-end for the detection task.
  \item \textbf{Generalization evidence:} the original paper found that YOLO degraded less than R-CNN on artwork datasets, suggesting useful global object representations.
\end{itemize}

\subsection{Limitations}
The same design choices that make YOLOv1 fast also impose important constraints:
\begin{itemize}
  \item Each cell predicts only two boxes and one shared class distribution, limiting the number of nearby objects it can represent.
  \item Several small or clustered objects with centers in the same cell are particularly difficult to detect.
  \item Repeated downsampling yields relatively coarse features for precise localization.
  \item The model can struggle with unusual aspect ratios, poses, and object configurations not well represented during training.
  \item The re-weighted SSE objective is a practical surrogate, but it is not perfectly aligned with AP or mAP, the final evaluation metrics.
\end{itemize}
The paper's error analysis identifies localization as YOLOv1's dominant source of error. These shortcomings provide useful context for understanding why later YOLO variants modify prediction heads, feature extraction, assignment strategies, and loss functions.

The limitations also explain why adding a second box predictor does not solve crowded-scene detection. If two nearby objects fall in one cell, their class supervision competes for one shared vector and the simple target assignment cannot represent an arbitrary set of instances. Increasing input detail alone cannot remove this output-capacity restriction. Meanwhile, the square-root size loss reduces one imbalance but still differs from directly optimizing box overlap: a small coordinate error can cause a large IoU decrease for a small object.

A geometric example makes this scale sensitivity concrete. Consider a square ground-truth box of side length $a$ and an equally sized prediction shifted horizontally by $d<a$. Their IoU is $(a-d)/(a+d)$. A shift of two pixels gives IoU $8/12\approx0.667$ for a ten-pixel square, but $98/102\approx0.961$ for a hundred-pixel square. Identical absolute coordinate errors can therefore have very different evaluation consequences. This example is a geometric calculation, not a measured failure from the network.

The confidence definition also creates a distinction between object presence and box quality. A low score can reflect poor expected overlap as well as weak evidence for the class. Increasing the score threshold may hide imprecise boxes, but can remove true detections rather than correcting localization. Conversely, a permissive threshold followed by NMS can reduce duplicates without repairing their coordinates. These mechanisms make qualitative inspection of boxes important alongside aggregate mAP; score filtering alone cannot overcome the representation and loss limitations described above.

\section{Typical Applications of YOLOv1}
\label{sec:applications}
The paper motivates fast object detection for responsive systems such as autonomous driving, assistive vision, robotics, and real-time webcam or video analysis \cite{redmon2016yolo}. A deployment decision should nevertheless account for the detector's known weaknesses in small-object and crowded-scene localization.

For a webcam demonstration or robot perception prototype, the appeal is a single forward pass that can update detections frequently. Video analysis can use the boxes as inputs to a separate tracking system, but YOLOv1 itself does not maintain identities across frames. Assistive perception benefits from low delay, yet distance estimation or action selection requires additional components beyond 2D boxes. Vehicle-scene analysis is another motivated application, although distant pedestrians and clustered road users directly expose the model's weaknesses.

These examples describe potential uses of the original detector, not claims that the paper validated every deployment. The VOC configuration covers 20 known categories and predicts axis-aligned boxes; it does not intrinsically provide instance masks, pose keypoints, unknown-category recognition, or temporal tracking. The suitability of an application therefore depends on its target classes, object scales, scene density, and acceptable localization error.

\section{Simple Demonstration / Implementation Example}
\label{sec:demo}
\noindent\textit{Reserved for later completion. The implementation has not yet been performed.}

\reservedfigure{Real-image inference example to be added.}{5.5cm}{Reserved space for input images and YOLOv1 inference results.}{fig:demo}

\FloatBarrier
\section{Evolution Beyond YOLOv1}
\label{sec:future}
\subsection{Reading the family as several research branches}
Later versions provide context for YOLOv1 rather than a single organizational lineage. Redmon and Farhadi developed v2/v3; later models came from several research teams. \Cref{tab:catalogue} records their sources. The progression in \cref{fig:evolution} summarizes technical themes rather than strict architectural inheritance.

\reportfigure{evolution_overview.pdf}{A concise map of technical themes beyond YOLOv1. Parallel branches pursue different combinations of architecture, training, and deployment improvements; the diagram summarizes the sources cited in this section.}{fig:evolution}

\subsection{From box priors to multi-scale representation: v2--v3}
YOLOv2 improves recall and localization with anchor boxes, dimension clustering, constrained center prediction, batch normalization, and Darknet-19. Passthrough features and multi-scale training improve flexibility. YOLO9000 additionally combines classification and detection supervision through hierarchical labels \cite{redmon2017yolo9000}.

YOLOv3 combines residual Darknet-53 with three prediction scales. Upsampling and concatenation preserve semantic and spatial information, helping small-object detection. Independent logistic predictions allow overlapping class labels. Weaker performance at stricter IoU thresholds shows that precise localization remains challenging \cite{redmon2018yolov3}. Together, v2/v3 address YOLOv1's coarse representation.

\subsection{Optimizing training and deployment: v4--v7}
YOLOv4 combines CSPDarknet53, spatial pyramid pooling, and path aggregation. Its distinction between bag-of-freebies and bag-of-specials separates training improvements from modifications that also affect inference. Mosaic augmentation and CIoU-based localization illustrate how better supervision can complement better features \cite{bochkovskiy2020yolov4}.

YOLOv5 develops a practical PyTorch pipeline with scaled model sizes. Later v5 revisions use C3 blocks and SPPF; AutoAnchor and training refinements support adaptation to a dataset. These revision-specific components should not all be attributed to the initial 2020 release \cite{ultralytics2020yolov5,yolov5architecture}. Here, v5 means the original anchor-based family, not the later anchor-free v5u variant.

YOLOv6 targets industrial efficiency using re-parameterizable feature extraction and an anchor-free decoupled head. Its v3.0 update adds bidirectional feature fusion, anchor-aided training, and localization distillation, keeping auxiliary training machinery out of the inference path \cite{li2022yolov6,li2023yolov6}. YOLOv7 instead emphasizes layer aggregation, architecture-aware scaling, planned re-parameterization, and lead-head-guided auxiliary supervision \cite{wang2023yolov7}. These parallel designs show that richer training need not require a correspondingly larger deployed graph.

\subsection{Revisiting the head, gradient, and post-processing: v8--v10}
YOLOv8's Ultralytics branch uses C2f blocks and an anchor-free split head. Task-aligned assignment and distribution-based box regression accompany its detection implementation \cite{ultralytics2023yolov8,ultralyticscode}. Removing predefined box shapes simplifies one modeling choice; it does not eliminate multiple predictions for one object. Its normal detection path still applies NMS.

YOLOv9 addresses unreliable learning signals caused by information loss through deep transformations. Programmable Gradient Information (PGI) supplies auxiliary supervision, while GELAN organizes efficient feature aggregation. The auxiliary structure supports training and can be removed for inference \cite{wang2024yolov9}. This extends concerns about gradient paths rather than merely adding more prediction layers.

YOLOv10 targets the cost of duplicate predictions and NMS. Consistent dual assignments combine one-to-many supervision with a one-to-one head; inference retains the latter and does not need NMS. Lightweight components and selective larger-context operations further optimize the accuracy--latency trade-off \cite{wang2024yolov10}. In contrast to YOLOv1's jointly learned prediction network, this use of end-to-end detection also removes NMS from the output pipeline.

\subsection{Toward attention-centric features: YOLO11--v12}
YOLO11 refines the Ultralytics line with C3k2 and C2PSA modules, adding selective attention while retaining a conventional anchor-free detection head \cite{ultralytics2024yolo11,ultralyticscode}. It does not automatically inherit YOLOv10's NMS-free assignment. Official software documentation supplies the references for YOLOv8 and YOLO11, since Ultralytics has not released formal original research papers for these models.

YOLOv12 makes efficient attention a central design concern through Area Attention and residual efficient layer aggregation (R-ELAN). Partitioning features into large areas reduces attention cost; residual aggregation supports stable optimization, and FlashAttention reduces memory-access overhead \cite{tian2025yolov12}. With a fixed number of areas, this reduces the coefficient of quadratic interaction cost rather than making attention linear in token count. Its benefits remain tied to implementation and hardware conditions.

The optimization scope expands from YOLOv1's output formulation to features, supervision, deployment, and context. Comparisons still require compatible AP definitions and latency conditions; the appendix catalogues contributions without ranking heterogeneous benchmarks.

\FloatBarrier
\section{Conclusion}
YOLOv1 demonstrates that object detection can be formulated as a single, global prediction problem. Its grid-based output, shared full-image features, and jointly trained detector provide a clear route to real-time inference. The same design restricts nearby-object capacity and precise localization, making the original paper useful both as a conceptual foundation and as an explanation of subsequent research needs. Later YOLO branches show how the common speed--accuracy objective can motivate changes to representation, supervision, post-processing, and attention rather than one universal sequence of architectural replacements. Implementation and qualitative inference results remain pending and will be added when supporting materials are available.

\clearpage
\appendix
\begin{landscape}
\section{YOLO Family Catalogue: v1--v12}
\label{sec:appendix}
The dates distinguish initial preprint/release from conference publication where relevant. The final column describes the detection problem or deployment emphasis; it is not a claim that each paper validated every downstream application. YOLO11 is the official Ultralytics name. All models listed concern fast object detection; some software families also provide separate variants for other tasks.

\begingroup
\small
\setlength{\tabcolsep}{5pt}
\renewcommand{\arraystretch}{1.17}
\begin{longtable}{@{}P{1.65cm}P{2.0cm}P{4.25cm}P{6.55cm}P{6.65cm}@{}}
\caption{Release history, responsible groups, critical improvements, and problem emphasis across the YOLO family.}\label{tab:catalogue}\\
\toprule
Version / source & First appearance & Individual / organization & Main characteristic or critical improvement & Target problem / task emphasis \\
\midrule
\endfirsthead
\multicolumn{5}{l}{\textit{Table \thetable\ continued}}\\
\toprule
Version / source & First appearance & Individual / organization & Main characteristic or critical improvement & Target problem / task emphasis \\
\midrule
\endhead
\midrule\multicolumn{5}{r}{\textit{Continued on the next page}}\\
\endfoot
\bottomrule
\endlastfoot
YOLOv1 \cite{redmon2016yolo} & 2015 preprint; 2016 CVPR & Joseph Redmon, Santosh Divvala, Ross Girshick, Ali Farhadi; University of Washington / AI2 / Facebook AI Research & Single full-image prediction; $7\times7$ grid; 24 conv + 2 FC; shared per-cell class vector & Replace proposal-based prediction with real-time, unified bounding-box detection. \\
YOLOv2 \cite{redmon2017yolo9000} & 2016 preprint; 2017 CVPR & Joseph Redmon and Ali Farhadi; University of Washington / AI2 / XNOR.ai & Darknet-19, learned anchor priors, stable center prediction, BN, multi-scale training; YOLO9000 joint supervision & Improve recall and localization; YOLO9000 additionally broadens category coverage using classification data. \\
YOLOv3 \cite{redmon2018yolov3} & 2018 & Joseph Redmon and Ali Farhadi; University of Washington & Residual Darknet-53; three prediction scales; independent logistic class outputs & Improve scale-dependent detection, particularly small objects, while retaining speed. \\
YOLOv4 \cite{bochkovskiy2020yolov4} & 2020 & Alexey Bochkovskiy, Chien-Yao Wang, Hong-Yuan Mark Liao; research collaboration including Academia Sinica & CSPDarknet53, SPP, PAN; coordinated BoF/BoS, Mosaic and CIoU & Improve the accuracy--speed balance with feasible single-GPU training. \\
YOLOv5 \cite{ultralytics2020yolov5,yolov5architecture} & 2020 repository & Ultralytics; led by Glenn Jocher & PyTorch pipeline; scaled sizes; original anchor-based head; later revisions use C3/SPPF and AutoAnchor & Make custom-dataset training, validation, export, and deployment practical. \\
YOLOv6 \cite{li2022yolov6,li2023yolov6} & 2022; v3.0 in 2023 & Chuyi Li and collaborators; Meituan & EfficientRep, re-parameterization, decoupled anchor-free head; v3.0 adds RepBi-PAN, AAT, DLD & Hardware-efficient industrial detection and deployment; retain richer supervision without auxiliary inference cost. \\
YOLOv7 \cite{wang2023yolov7} & 2022 preprint; 2023 CVPR & Chien-Yao Wang, Alexey Bochkovskiy, Hong-Yuan Mark Liao; Academia Sinica collaboration & E-ELAN, concatenation-aware scaling, planned re-parameterization, trainable bag-of-freebies & Improve feature learning and speed--accuracy scaling through architecture and auxiliary training. \\
YOLOv8 \cite{ultralytics2023yolov8,ultralyticscode} & January 2023 & Ultralytics; Glenn Jocher, Ayush Chaurasia, Jing Qiu and contributors & C2f, anchor-free split head, distributional regression; unified API & Simplify detection training/deployment; separate variants support segmentation, pose, classification, and OBB. \\
YOLOv9 \cite{wang2024yolov9} & 2024 & Chien-Yao Wang, I-Hau Yeh, Hong-Yuan Mark Liao; Academia Sinica / National Taipei University of Technology & Programmable Gradient Information and GELAN & Address information loss and improve reliable learning signals in efficient detection networks. \\
YOLOv10 \cite{wang2024yolov10} & 2024 & Ao Wang and collaborators; Tsinghua University-led research team & Consistent dual assignments, one-to-one inference, NMS-free output; efficiency-driven component design & Reduce duplicate prediction and post-processing latency for end-to-end real-time detection. \\
YOLO11 \cite{ultralytics2024yolo11,ultralyticscode} & September 2024 & Ultralytics; Glenn Jocher, Jing Qiu and contributors & C3k2, C2PSA, refined feature extraction; conventional anchor-free detection head & Improve efficient detection in the Ultralytics ecosystem; distinct variants cover multiple vision tasks. \\
YOLOv12 \cite{tian2025yolov12} & February 2025 & Yunjie Tian, Qixiang Ye, David Doermann; University at Buffalo / University of Chinese Academy of Sciences & Area Attention, R-ELAN, attention-oriented optimization, FlashAttention & Capture broader feature interactions while constraining latency and stabilizing attention-based training. \\
\end{longtable}
\endgroup
\end{landscape}

\clearpage

\begingroup
\small\setstretch{1.0}
\bibliographystyle{unsrt}
\bibliography{references}
\endgroup

\end{document}
